\chapter{Conclusion and Future Work}
\label{chapter:conclusion}

\section{Summary}

This thesis set out to establish water-skipping as a usable mode of
locomotion for a micro aerial vehicle, and to show that a small rotorcraft can
be thrown back into the air by its own spinning hydrofoils.

The mechanism rests on a single observation: a surface that is nearly useless
in air becomes highly effective the instant it enters water, because the two
fluids differ in density by three orders of magnitude.  A vehicle that spins
continuously can therefore carry hydrofoils at little cost to its flight, and
call on them for a large impulse during a contact lasting only tens of
milliseconds.  Chapter~\ref{dynamic} developed a quasi-steady blade-element
model of that contact, which identifies the inclination of the foils and the
conditions of entry as the quantities that govern the outcome, and shows the
central tension of the manoeuvre: the same sweep that generates the ejecting
force is itself consumed by the water.

Realising this required a vehicle whose spin could be set rather than merely
tolerated.  Chapter~\ref{chapter:paperone} describes an airframe in which two
rotors are mounted tangentially and drive the rotation directly, while two
remain vertical and provide lift and attitude authority.  Decoupling the yaw
drive from the lift channel makes the spin rate an operating parameter chosen
by the experimenter, and, more consequentially for this work, means the
rotation persists when the lift command is removed entirely.  The flight
control that keeps such a vehicle stable and position-controllable while
rotating is developed in Chapter~\ref{chapter:papertwo}, together with the
online separation of the spin-synchronous artefact from the attitude estimate
that makes closed-loop control practical at these rates.

Chapter~\ref{chapter:paperthree} reports the flight experiments.  The vehicle
descends onto the water with its lift rotors commanded to zero, and is
returned to the air by the hydrofoils alone; the controller is re-engaged only
once the vehicle is already climbing away from the surface.  Hops were
repeated, with the vehicle returning to its commanded height between contacts,
demonstrating that the full cycle closes rather than that a single impact can
be survived.

\section{Contributions}

The principal contribution is the locomotion mode itself: a rotorcraft that
travels across water in a sequence of hops, using hydrofoils driven by its own
rotation to recover a substantial part of the energy of each contact.  This
treats the water surface as something to move across rather than as a place to
land, dive, or sample.

Supporting this are a quasi-steady model of hydrofoil-driven ejection that
relates the hop to foil geometry and entry conditions; an airframe whose spin
rate is commanded independently of its lift, which both makes the spin a
design variable and allows the descent to be flown unpowered; and flight
control at spin rates approaching
$10{,}000^{\circ}\,\mathrm{s^{-1}}$, at which no controllability limit was
encountered.

\section{Limitations}

\paragraph{The model over-predicts the contact.}
The blade-element treatment of Chapter~\ref{dynamic} assumes that the foils
are immersed to their geometric depth in an undisturbed, flat surface.  At the
sweep speeds this vehicle reaches, that assumption does not hold: the measured
contact is markedly longer and gentler than predicted, which is consistent
with the free surface deforming away from the foils and with ventilation of
the upper surface.  The model remains useful for the trends it exposes and for
sizing, but its absolute predictions should not be relied upon at these
speeds.

\paragraph{Inertia was never measured.}
The moment of inertia about the spin axis is estimated from a uniform-disk
approximation rather than measured.  Predicted spin loss through the contact
depends on it directly, and should be treated as indicative until it is
determined experimentally.

\paragraph{Instrumentation limited the length of a sequence.}
The number of consecutive hops that could be recorded was limited by the
motion-capture system rather than by the vehicle.  Water thrown up at contact
intermittently interrupts marker tracking at the moment the controller most
needs position feedback.  Mitigations reduced the cost of a dropout to a
single hop, but the limitation is one of the measurement environment, and an
onboard state estimate would remove it.

\section{Future Work}

\paragraph{Contact hydrodynamics.}
The clearest opening is the discrepancy identified above.  Modelling the
deformation of the free surface during entry, rather than assuming it flat,
would close the gap between the predicted and measured contact and turn the
model into a quantitative design tool.  This is also the route to predicting
the depth to which the foils penetrate, which the present treatment cannot.

\paragraph{Disturbed and moving water.}
Every hop reported here was made onto a still surface.  Water encountered
outside a laboratory is not still, and the surface state at the moment of
contact is an untouched variable: waves change the effective entry angle, the
local surface velocity, and the depth available to the foils.  Whether the
vehicle can hop reliably into a disturbed surface, and how the foils should be
shaped if it must, is a natural continuation of this work.

\paragraph{Onboard estimation.}
Removing the dependence on external motion capture would lift the constraint
that limited hop sequences here, and is a precondition for operating outside a
prepared volume.

\paragraph{Phase lag at high spin rate.}
No controllability limit was reached in this work, and the spin rates flown
were well above what earlier reasoning suggested would be workable.  A
plausible explanation is that the binding constraint is the phase lag around
the closed loop rather than the number of control updates per revolution,
since a rotating vehicle turns measurably between a measurement being taken
and the corresponding command taking effect.  This was not investigated: no
latency measurements were recorded, and the observation is offered only as a
direction, not as a result.

\paragraph{Foil design.}
Only one hydrofoil geometry was fabricated and flown.  The inclination and
planform were selected from the model rather than from experiment, and a study
that varies them on hardware would test the model's predictions directly and
is the obvious follow-on to the present work.
